Vor dem Eintauchen zeigt der Kraftmesser die Gewichtskraft $F_\text{G} = \rho_\text{Al}\,V\,g$ an, danach $F = F_\text{G} - F_\text{A}$. Der Auftrieb ist also
\begin{align}
 F_\text{A} &= \FAkF = \FAk
\end{align}
Der Auftrieb ist die Gewichtskraft der verdrängten Flüssigkeit, $F_\text{A} = \rho\,V\,g$. Teilt man die beiden Kräfte durcheinander, fallen $V$ und $g$ heraus:
\begin{align}
 \frac{F_\text{A}}{F_\text{G}} &= \frac{\rho}{\rho_\text{Al}}
\end{align}
Mit $\rho_\text{Al} = \rhoAlO$ (FoTaBe S.~187) folgt
\begin{align}
 \rho &= \rhoFlF \\
   &= \frac{\FG-\FT}{\FG}\cdot\rhoAl \\
   &= \rhoFl \approx \result{\rhoFlP{0}{2}}
\end{align}
Das ist etwa die Dichte eines Öls.
